An operator path describes the forcing, but we still need to identify the wave
it drives. There is a small but important issue with the initial data: the
constant displacement is not in the whole-space Sobolev phase used for the
operator estimates. We therefore keep the constant coordinates explicitly while
retaining both displacement and velocity.

The input is the path construction and source convergence from
Chapter~\ref{S:book-operator-paths}, together with the canonical expansion from
Chapter~\ref{S:G0}. Section~\ref{S:masked-moments} compares the regularized
solutions and their tested moments. Section~\ref{S:constant-phase-young}
constructs Young evolution on the enlarged phase space and identifies it with
the physical wave. Finally, Section~\ref{S:compact-mean-family} studies compact
families of nonnegative deterministic potentials. Their positive ground states
and uniform spectral gap will control growth in the lower-bound argument. All
masks and horizons here are fixed. These results do not yet justify letting the
spatial radius grow with the observation time.
